% !TeX root = ../教材题库_全部条目_转写工作稿.tex
% Source: XJ-B1-C05-TH-P163-Y0689
% 第5章 / 5.2 任意角的三角函数 / 教材 p.163、164 / PDF 第170、171页
% OCR未校对：公式与文字均可能有误；未完成数学TeX转写；禁止直接发给学生
% 原题及相邻原解截图（核对用）：../题目原图/XJ-B1-C05-TH-P163-Y0689-01.webp, ../题目原图/XJ-B1-C05-TH-P163-Y0689-02.webp
% 本文件只在转写完成后作为题干使用；OCR 条目尚不是可靠题干。
\begin{bankitem}{p.163 思考（含紧邻说明）}{第5章 / 5.2 任意角的三角函数 / 教材 p.163、164 / PDF 第170、171页}{XJ-B1-C05-TH-P163-Y0689}{OCR未校对：公式与文字均可能有误；未完成数学TeX转写；禁止直接发给学生}
\begin{ocrdraft}
[教材p.163 / PDF第170页]
\par
？思考任意给定一个角仅，以角仅的顶点为原点0，以它的始边为x轴的非负
\par
半轴，建立平面直角坐标系．在角终边OM上任取不同于原点0的点P(x，回，
\par
记OP=丆=2十，则，还和的值都分别是唯一确定的值吗？
\par
若角的终边落在第一象限，如图5，2一1所示．过点P
\par
作工轴的垂线，垂足为点D.于是，OD=x，DP一3，，OP=
\par
7一十．在终边OM上另取一点p气，），过点p'作
\par
轴的垂线，垂足为点0．于是，00一x'，D'P'一y',
\par
0
\par
工2十3，2．由直角三角形和相似三角形知识可得
\par
图5．2一1
\par
这表明：任意给定一个在第一象限的角仅，相应的六还和的值都分别是唯一
\par
确定的值，且与终边上的点P的位置无关．
\par
类似地，我们可以讨论出，当角a的终边分别落在第二
\par
象限、第三象限、第四象限、x轴上时，上述结论仍然成立．
\par
若角a的终边落在y轴上，则除了无意义外，上述结
\par
论仍成立
\par
0
\par
一般地，如图5．2一2，任意给定一个角a，以角仅的顶
\par
图5．2一2
\par
点为原点0，以它的始边为工轴的非负半轴，建立平面直角
\par
坐标系，则对于角终边OM上任一点P(x，回，
\par
，莖和（终边不在3，轴上）的值
\par
[教材p.164 / PDF第171页]
\par
都分别是唯一确定的，其中厂一十重．
\par
卫
\par
工的值分别称为角a的正弦、余弦、正切，且分
\par
数学上，将这样得出的
\par
别记作sina，cosa，tan伍即
\par
一，其中厂一十3，2．
\par
Sln仅．一
\par
COS仅．一
\par
tan
\end{ocrdraft}
\end{bankitem}
